\documentclass{article}
\usepackage{amsmath}
\usepackage{array}

\begin{document}

\section*{Collocation Method for a Second-Order ODE}

\subsection*{1. The Initial Problem}
We seek a solution $u(t)$ for the following boundary value problem (BVP) on the interval $t \in [0, 1]$:
\begin{equation}
    u''(t) = 0
\end{equation}
subject to the boundary conditions:
\begin{itemize}
    \item $u(0) = 0$
    \item $u(1) = 1$
\end{itemize}
and the interior collocation requirements:
\begin{itemize}
    \item $u''(1/3) = 0$
    \item $u''(2/3) = 0$
\end{itemize}

\subsection*{2. Case A: Specific Power Basis}
Using the standard monomial basis $\phi = \{1, t, t^2, t^3\}$, where $u(t) = a_0 + a_1t + a_2t^2 + a_3t^3$, the system $Ax=b$ is constructed as:
\begin{equation}
\begin{bmatrix} 
1 & 0 & 0 & 0 \\
1 & 1 & 1 & 1 \\
0 & 0 & 2 & 2 \\
0 & 0 & 2 & 4
\end{bmatrix}
\begin{bmatrix} 
a_0 \\
a_1 \\
a_2 \\
a_3
\end{bmatrix} = 
\begin{bmatrix} 
0 \\
1 \\
0 \\
0
\end{bmatrix}
\end{equation}

The determinant of the matrix $A$ is $4$, indicating a unique solution exists.

The solution to which is:
\begin{equation}
\begin{bmatrix} 
a_0 \\ a_1 \\ a_2 \\ a_3
\end{bmatrix} = 
\begin{bmatrix} 
0 \\ 1 \\ 0 \\ 0
\end{bmatrix}
\end{equation}

The matrix is decomposable into a linear combination of pauli strings.
In this case we need 16 out of 16 total pauli strings to represent the matrix exactly.
% does interchancing rows or colimns change the number of pauli strings needed? Seems like not.


\subsection*{3. Case B: Generalized Basis with Coefficients $C$}
We now define the basis functions $\phi_i(t)$ as general polynomials with coefficients $C_{ij}$:
\begin{align*}
    \phi_0(t) &= C_{00} \\
    \phi_1(t) &= C_{11}t + C_{10} \\
    \phi_2(t) &= C_{22}t^2 + C_{21}t + C_{20} \\
    \phi_3(t) &= C_{33}t^3 + C_{32}t^2 + C_{31}t + C_{30}
\end{align*}
Applying the four conditions to the trial function $u(t) = \sum_{i=0}^3 a_i \phi_i(t)$ yields the system:
\begin{equation}
\begin{bmatrix} 
C_{00} & C_{10} & C_{20} & C_{30} \\ 
C_{00} & (C_{11} + C_{10}) & \sum C_{2j} & \sum C_{3j} \\ 
0 & 0 & 2C_{22} & (2C_{33} + 2C_{32}) \\ 
0 & 0 & 2C_{22} & (4C_{33} + 2C_{32}) 
\end{bmatrix}
\begin{bmatrix} 
a_0 \\ 
a_1 \\ 
a_2 \\ 
a_3 
\end{bmatrix} 
= 
\begin{bmatrix} 
0 \\ 
1 \\ 
0 \\ 
0 
\end{bmatrix}
\end{equation}

The determinant of this matrix is given by:
\begin{equation}
4 C_{00} C_{11} C_{22} C_{33}
\end{equation}

% in general, is the determinant always a product of the leading coefficients?

Decomposing the matrix into the sum pauli strings, we find that the number of required pauli strings varies with the choice of coefficients $C_{ij}$.
One can see that the minimum number of pauli strings needed is 3, for example with the choice:
\begin{equation}
    -3 II - 2 IX + IZ
\end{equation}

Which gives the Coefficients

The matrix decomposition into Pauli strings reveals interesting structure. For the choice of coefficients:
\begin{align*}
    C_{00} &= 2, \\ 
    C_{10} &= 2, \quad C_{11} = 2 \\
    C_{20} &= 0, \quad C_{21} = -1, \quad C_{22} = 1 \\
    C_{30} &= 0, \quad C_{31} = -1, \quad C_{32} = 0, \quad C_{33} = 1
\end{align*}


Which is a 81\% reduction from the 16 pauli strings needed in Case A.



\end{document}



% explain how you got the pauli strings
% run some code numerically justify. Run for two cases. The case B should look better. Add some figures

